\begin{proof}

\begin{enumerate}[label={$\langle1\rangle$\arabic*.},leftmargin=*]
\item On a marked smooth torus, the exponential sequence presents $\Pic^0$ as constant complex lattice cocycles modulo the integral cocycles and the two holomorphic period rows. The equation for $z\gamma$ to belong to this sum is
\[
 z\gamma=(n,p,q,r)+a(6\mu,\tau,1,0)+b(\beta_c,\mu,0,1).
\]
The last two coordinates give $a=-q,b=-r$, and the first two give \eqref{cdp:eq:C.10}. If $r=0$, positivity of $\operatorname{Im}\tau$ forces $p=q=0$ and $z=n$. Otherwise
\[
 \operatorname{Im}z=-r\left(\operatorname{Im}\beta_c-
           6(\operatorname{Im}\mu)^2/\operatorname{Im}\tau\right)\ne0.
\]
In particular real nonintegral twists have no regular resonance. The fixed values $\mu(z_1)=(2-\rho)/3$ and $\mu(z_2)=(1-i)/2$ give the two displayed lattices by solving the first equation. Their nonreal generators make each representation $n+k\lambda_j$ unique.

\item We give a coherent local complex to calculate every length. On a smooth simply connected torus-family disc use normalized period coordinates $\zeta=Z(s)x+y$, $x,y\in\R^2/\Z^2$, and translation-invariant forms $G_j=(dx_j)^{0,1}$. The invariant horizontal lift is $h=\partial_s+\sum_i(Z'(s)x)_i\partial_{\zeta_i}$. Since $dx_j$ is closed, $\bar\partial G_j=0$ and
$\bar\partial h=\sum Z'_{ij}e_iG_j$. After subtracting the integer cocycle of a resonance, the line multiplier has invariant Dolbeault form $a_1G_1+a_2G_2$, up to the common nonzero factor $2\pi i$. Thus the invariant relative complex has terms
\[
 C^j=(V\oplus\OO h)\otimes\Lambda^jG,\qquad j=0,1,2,
\]
\[
 d(v\eta)=v(a\wedge\eta),\qquad
 d(h\eta)=h(a\wedge\eta)+\kappa\wedge\eta,             \tag{C.13}\hypertarget{tag-008R}{}\label{cdp:eq:C.13}
\]
where
\[
 a=(z-n+6q\mu+r\beta_c,-p+q\tau+r\mu),\qquad
 \kappa=\begin{pmatrix}6\mu'&\tau'\\\beta_c'&\mu'\end{pmatrix}.
\]
This complex computes the derived direct image. To verify the comparison, decompose relative smooth forms into Fourier characters under real torus translations. The zero character gives \eqref{cdp:eq:C.13}. For every other character the relative Dolbeault differential is exterior multiplication by a nonzero covector; contraction by its normalized dual is a homotopy. On a sufficiently small closed subdisc those covectors are bounded away from zero at bounded frequency and grow linearly at large frequency. The homotopies preserve smooth Fourier series and all their derivatives. The vertical-horizontal filtration of the tangent extension, followed by the filtration in the base antiholomorphic differential, gives a finite filtration with acyclic nonzero-character quotients. Hence their full complexes are acyclic. The zero-character formulas have holomorphic coefficients, so the remaining base Dolbeault resolution is precisely the resolution of \eqref{cdp:eq:C.13}. A finite cyclic stabilizer normalizes translations, so this construction and its averaging projection are equivariant on each root chart.

\item At a regular resonance, $a'=\kappa^t(q,r)^t$ and $(q,r)\ne(0,0)$. The period determinant calculation in Part E shows that $\kappa$ is invertible at every regular point. Thus $a$ has a simple zero; choosing its direction as the second cocycle frame and changing the vertical frame gives
\[
 a=(0,s),\qquad\kappa=\begin{pmatrix}0&1\\1&0\end{pmatrix}.
\]
In degree one the two generators $e_1G_2,e_2G_2$ are killed by $s$. In degree two the relations eliminate the two vertical generators and leave $hG_1G_2$ killed by $s^2$. This gives the regular modules in \eqref{cdp:eq:C.12}.

\item At the order-three point put $t=s^3$, $\nu=\mu-(2-\tau)/3$ and $b=\beta_c+4\mu+(2/3)\tau-b(0)$. The order-three characters give $\nu=\nu_2s^2+O(s^3)$ and $b=b_3s^3+O(s^4)$. We check the joint coefficient that controls the module. With elliptic scale $f(s)=Fs^2+O(s^3)$, the Weierstrass equation gives
\[
 \nu_2^2=9F^2/512,\quad \tau_1=i\pi F^2/2,\quad
 6\mu'^2-\tau'\beta_c'=(3F^2/4)s^2+O(s^3).
\]
The first equality follows from the exact tripling coordinate $x(3P)(0)=512/81$ and $x=f^2\wp$: $3\nu$ approaches zero, so $x(3P)(0)=F^2/(9\nu_2^2)$. The determinant equality follows from the identity in Part E and $f^4g_2=12t^3(t-1)$, $f^6g_3=-8t^4(t-1)^2$. Comparing coefficients in $6\nu'^2-\tau'b'$ gives $b_3=7i/(32\pi)$. For a nonintegral resonance the Picard section is $3k(\nu,b)$ in the corresponding frames. Eliminating the unit $\tau'$ leaves transverse coefficient
\[
 3k(b-6\nu\nu'/\tau')
   =\frac{123ik}{64\pi}s^3+O(s^4),\qquad k\ne0.        \tag{C.14}\hypertarget{tag-008S}{}\label{cdp:eq:C.14}
\]
It is nonzero. Equivariant holomorphic frame operations therefore reduce \eqref{cdp:eq:C.13} to
\[
 a=(0,s^2),\qquad\kappa=\begin{pmatrix}0&s\\s&1\end{pmatrix}.
\]
Its degree-one and degree-two modules are both $\C\{s\}/(s^2)\oplus\C\{s\}/(s^3)$. Explicit degree-one cycles are $u=e_1G_2$, $v=e_2G_1+shG_2$; the third cycle is removed by the boundary $su+e_2G_2+sv$. The remaining relations are $s^2u=s^3v=0$. In degree two elimination of $e_2G_1G_2$ leaves the generators $e_1G_1G_2,hG_1G_2$ with annihilators $s^2,s^3$.

\item We retain the actual affine characters. The integral cocycle at the order-three resonance is $(n,-2k,-k,-3k)$; evaluating on $v_1$ gives $n$. At the order-four resonance it is $(n,-k,-k,-2k)$; evaluation on $v_2$ gives $-n\equiv3n\pmod4$. The invariant exponents at the order-three point are $n+2,n+2$ for the two degree-one generators and $n+2,n$ in degree two. Here an invariant exponent is the required Taylor exponent modulo the stabilizer order. Thus the length-two summand retains precisely the monomials $s^j$ with $0\le j<2$ and $j\equiv n+2\pmod3$: its invariant lengths for $n\equiv0,1,2$ are $0,1,1$. A length-three summand retains one coefficient in every residue class. Both coarse lengths are therefore $1+\epsilon_n$, and multiplication by $t=s^3$ kills them. At the order-four point the regular normal form applies. The two length-one degree-one generators have invariant exponents $3n+1,3n+2$, and the length-two degree-two generator exponent $3n+2$. Their invariant length is one exactly for $n\equiv1,2\pmod4$, and the coarse parameter $s^4$ kills them. This proves all noncusp summands.

\item At the cusp, the restriction of $M_z$ has normalization pullback trivial and dual-complex character $(a,1)$, $a=e^{2\pi iz}$. Its normalization complex is
\[
 \C\xrightarrow{(a-1,0)^t}\C^2
 \xrightarrow{(0,a-1)}\C.
\]
For $z\notin\Z$ it is acyclic. Semicontinuity and base change make $R\bar f_*M_z$ vanish on a cusp neighbourhood. The logarithmic sequence (C.8) then gives $R\bar f_*(\mathcal T\otimes M_z)=0$ there. Sequence (C.3) contributes exactly the reduced skyscraper $\C_{p_0}$ in degree two for the ordinary tangent bundle. Away from resonances a nontrivial topologically trivial line on a torus has zero cohomology: its unitary Fourier presentation has no zero eigenmode; the tangent extension by trivial bundles retains that acyclicity. We verify that this analytic support is proper. Fix a simply connected disc in the regular base. For each fixed integer tuple $(n,p,q,r)$, a resonance has $r\ne0$ and $a^\prime=\kappa^t(q,r)^t\ne0$, so its zeros are isolated. Each such zero set is countable on this disc, and the union over the countably many tuples is countable. The disc has a point outside this union. At that fibre the preceding line-bundle acyclicity and tangent extension give zero cohomology in every degree. Upper semicontinuity gives a neighbourhood with the same vanishing. Each coherent higher direct image therefore has proper analytic support on the compact base curve, hence finite support. Every regular resonance lies in that support by its computed nonzero local modules, proving that $\mathcal R_z$ is finite. Degree zero is zero because a direct image of a bundle is torsion-free and has generic rank zero. These arguments prove \eqref{cdp:eq:C.12}, and Leray proves \eqref{cdp:eq:C.11}. The logarithmic degree-two direct image is a torsion sheaf for these twists, so its first cohomology vanishes. Equation (C.4) proves the final surjectivity assertion.

\end{enumerate}
\end{proof}
